\chapter{Summary}\label{sec:summary}

This work examined classical and neural network–based relay strategies for \textit{two-hop relay communication} using a unified simulation framework. The objective was to analyze performance–complexity tradeoffs and determine when learning-based relays provide practical advantages over analytically derived methods.
The core comparison is carried out on one canonical setup, never departed from within it: a two-hop SISO link over i.i.d.\ Rayleigh fast fading, complex baseband, QPSK, uncoded BER, with the relay function as the only varied axis. The coded study and the unknown-channel contribution depart from it deliberately. The study showed the following:

\begin{enumerate}
    \item The best tested learned relays outperform AF at the evaluated low-SNR points (H1); the experiments do not isolate a unique mechanism for that advantage.
    \item Classical \textit{decode-and-forward (DF)} matched or outperformed all neural approaches at \textbf{medium-to-high SNR} while requiring no trainable parameters, confirming that DF is the best-performing evaluated symbol-wise classical relay in this regime (H2).
    \item Increasing model size did not produce a \textbf{systematic BER improvement}, and the replicated size sweep found no statistically meaningful upward BER trend attributable to model capacity (H3).
    \item The 169-parameter MLP is competitive with larger models in the original canonical comparison. This is not a universal minimum-size result, and no train/test-gap study establishes overfitting as the explanation.
    \item The replicated size sweep identifies the smallest tested configurations meeting each degradation criterion: four parameters suffice for three memoryless cases at zero degradation, while the selected three-tap cases use 73--145. Window length matters in this grid, but the experiment does not prove architecture lower bounds.
    \item At an equal parameter budget, feedforward and sequence relays remained separated at low SNR. The apparent high-SNR convergence was smaller than the Transformer run-to-run spread and is not seed-robust; equalizing parameters also left depth, width, state dimension, window, and optimization confounded (H4).
    \item Within the simulated benchmark, a \texttt{hybrid relay} combining neural processing at low SNR with DF at high SNR tracked the best evaluated relay when operating SNR was known. Its switching-threshold robustness, quantization, energy use, synchronization, and hardware behavior were not evaluated.
\end{enumerate}
    Findings 1--4 and 6--7 concern the canonical comparison; the size sweep spans several specified channel families. The following items summarize Chapter~\ref{sec:unknown-channels}. The supplementary coded study is summarized below.
\begin{enumerate}
\setcounter{enumi}{7}
    \item On the selected BPSK three-tap filter, AF and zero-threshold DF approach the derived $0.25$ floor, while the 170-parameter MLP substantially reduces BER through the fixed-budget 16-dB point. Matched genie MLSE leads over the resolved range; the interpolated required-SNR penalties are $0.26$, $1.18$, $1.88$, and $2.53$~dB at BER $10^{-1}$ through $10^{-4}$. These are not a constant offset. In the separate blind study, the family-trained MLP narrowly leads the evaluated CMA implementation without test-time adaptation.

    \item Sweeping the intermediate \emph{partial}-posterior case at 10~dB located a crossover between 20 and 10 pilots; its SNR dependence was not measured. Pilot-aided LS-plus-MLSE wins when its estimate is sufficiently accurate, while the family-trained relay avoids pilot overhead and per-block estimation. The result is specific to the tested estimator, channel family, block lengths, and SNR.

    \item The separate coded-QPSK memory sweep uses a window-11, eight-hidden-unit, four-output MLP with 220 parameters and 208 MACs per symbol. The nominal full-state MLSE count is $2M^L$ under the stated convention; the two counts cross between the evaluated three- and five-tap cases and differ by about $158\times$ at seven taps. MLSE remains more accurate. Fixed arithmetic is a property of this deployed network, not a guarantee that its accuracy or required size remains fixed for longer channel memories.

    \item \textbf{Latency did not separate the learned relay from sequence detection at the available resolution} on the tested channel. No BER difference was resolved among the tested traceback depths; a three-symbol MLSE setting had fewer symbols of decision delay than the window-11 relay's five. The block relay required a full frame before producing output and showed no resolved accuracy gain in this experiment. The two research objectives of Chapter~\ref{sec:research-objectives} therefore have boundaries set by different resources in the comparisons evaluated here --- channel-estimate reliability for the unconstrained question, arithmetic for the constrained one --- and the distance between them is a result rather than a restatement.
\end{enumerate}

\begin{longtable}[]{@{}
  >{\raggedright\arraybackslash}p{(\linewidth - 4\tabcolsep) * \real{0.3333}}
  >{\raggedright\arraybackslash}p{(\linewidth - 4\tabcolsep) * \real{0.3333}}
  >{\raggedright\arraybackslash}p{(\linewidth - 4\tabcolsep) * \real{0.3333}}@{}}
\caption[Hypothesis outcomes summary]{Research hypotheses, their statements, and experimental outcomes}\label{tbl:table33}\tabularnewline
\toprule\noalign{}
\begin{minipage}[b]{\linewidth}\raggedright
Hypothesis
\end{minipage} & \begin{minipage}[b]{\linewidth}\raggedright
Statement
\end{minipage} & \begin{minipage}[b]{\linewidth}\raggedright
Result
\end{minipage} \\
\midrule\noalign{}
\endfirsthead
\toprule\noalign{}
\begin{minipage}[b]{\linewidth}\raggedright
Hypothesis
\end{minipage} & \begin{minipage}[b]{\linewidth}\raggedright
Statement
\end{minipage} & \begin{minipage}[b]{\linewidth}\raggedright
Result
\end{minipage} \\
\midrule\noalign{}
\endhead
\bottomrule\noalign{}
\endlastfoot
H1 & Some AI relays outperform AF at low SNR & \textbf{Confirmed}: the best neural relays beat AF at 0--4 dB on the canonical channel \\
H2 & DF lowest-BER among per-symbol relays at high SNR & \textbf{Confirmed}: DF matches or beats all AI methods at \(\geq 6\) dB with 0 parameters \\
H3 & An error minimum at intermediate capacity & \textbf{Not supported in the tested range}. The replicated size sweep finds no robust BER upturn with capacity. No train/test-gap analysis establishes an overfitting mechanism; uncertainty is configuration-dependent \\
H4 & Equal parameter count makes architecture-dependent differences negligible & \textbf{Not supported across SNR}. Low-SNR differences remain; the apparent high-SNR convergence is not robust to all measured seeds. Width, depth, windows, state dimensions and optimization remain confounded \\
H5 & A small learned relay improves on AF and zero-threshold DF under the selected model mismatches & \textbf{Supported within the evaluated configurations}. The BPSK fixed-budget MLP reaches $1.20\times10^{-7}$ at 16~dB (36 errors in 300M bits), but does not establish an advantage over matched MLSE. QPSK has a substantial matched-MLSE advantage over the MLP. Pilot and blind results concern the specified impairment family, classical implementations, and budgets \\
\end{longtable}

The evidence supports targeted use of learned relays, not a general learned-versus-classical verdict. The measured advantages depend on model mismatch, available channel information, architecture, and budget. Reduced-state sequence estimation, decision-feedback equalization, and fully reliability-weighted decoding remain unmeasured alternatives.

The supplementary coded study reports a high-SNR reversal between learned forwarding and the implemented block-DF pipeline. Its decoder metrics ignore symbol-dependent post-equalization reliability, so the observations do not establish that decoding at a relay inherently wastes redundancy or must lose at high SNR. The rate-adaptation experiment compares two classical forwarding schemes on a finite modulation-and-coding grid; its goodput ordering and idealized buffering constraints are not a neural-relay deployment guarantee. Hardware timing, full half-duplex deadline accounting, other code families, and generalization outside the trained impairment families remain beyond the experiments.
